\chapter{Esperimento con gli utenti}
\label{ch:esperimento}

Una volta conclusa la riprogettazione del sito sono stati eseguiti degli esperimenti con gli utenti: lo scopo di questa fase è stato quello di raccogliere dati relativi all'utilizzo del sistema da parte di utenti reali per poi usarli, nel capitolo successivo, per analizzare nel dettaglio gli effettivi miglioramenti apportati.
Il test ha confrontato le due versioni del sistema --- V1 (originale) e V2 (riprogettata) --- su un campione temporale di partecipanti. In questo capitolo verranno illustrate le modalità con cui sono stati svolti questi esperimenti.

\section{Caratteristiche dell'esperimento}

Per gli esperimenti di usabilità è stato scelto l'approccio "within-subject", che implica che ogni gruppo di utenti abbia testato entrambe le condizioni sperimentali, ossia le due versioni del sito, in ordine variabile: questa metodologia è preferibile in contesti nei quali sono ridotti i problemi di apprendimento, richiede meno utenti e comporta minori problemi legati alla variabilità degli stessi.

È stata inoltre utilizzata la metodologia di \textit{counterbalancing} per evitare problemi legati all'apprendimento degli utenti da una versione all'altra. Un totale di 10 utenti è stato selezionato e diviso in due gruppi distinti:
\begin{itemize}
    \item \textbf{Gruppo A}: ha sperimentato inizialmente la "Versione 1" del sito (originale) seguita dalla "Versione 2" (riprogettata).
    \item \textbf{Gruppo B}: ha sperimentato inizialmente la "Versione 2" del sito seguita dalla "Versione 1".
\end{itemize}

Si è cercato, per quanto possibile, di garantire omogeneità e bilanciamento tra i gruppi considerando le informazioni demografiche raccolte attraverso questionari preliminari (Sesso, Età, Titolo di studio, Conoscenze informatiche).

\section{Ambiente}

Per simulare il contesto di utilizzo reale del sistema, si è deciso di eseguire gli esperimenti in luoghi familiari agli utenti (come aule studio e abitazioni private), in modo da metterli a loro agio.
Tutti gli esperimenti sono stati eseguiti su un computer portatile fornito dallo sperimentatore (con browser Chrome), in modo da ricreare la stessa situazione e ambiente applicativo (Laravel locale) per tutti i soggetti.
Lo sperimentatore si trovava nella stessa stanza in posizione non intrusiva, intervenendo esclusivamente su esplicita richiesta del partecipante.

\section{Dati raccolti}

I dati raccolti dallo sperimentatore sono sia quantitativi che qualitativi, così da permettere sia analisi statistiche che valutazioni di natura descrittiva.
Per ogni compito sono stati raccolti i seguenti dati quantitativi:
\begin{itemize}
    \item \textbf{Completamento} (indica se il task è stato completato, completato con aiuto, o non completato);
    \item \textbf{Tempo impiegato} per il completamento del compito;
    \item \textbf{Numero di click} necessari al completamento;
    \item \textbf{Numero di errori} eseguiti durante la navigazione.
\end{itemize}

Per i dati qualitativi si hanno:
\begin{itemize}
    \item \textbf{Note sperimentatore};
    \item \textbf{Commenti dell'utente} (\textit{Think Aloud}).
\end{itemize}

Inoltre ogni utente ha valutato la \textbf{difficoltà} di ogni scenario su una scala appropriata e ha compilato il questionario \textbf{SUS (System Usability Scale)} per il grado di usabilità, affiancato dall'indice \textbf{NPS (Net Promoter Score)}.

\section{Scelta degli sperimentatori}

Per l'esecuzione degli esperimenti si è scelto di impiegare un solo sperimentatore, responsabile dell'osservazione e dell'annotazione delle metriche, con funzione anche di facilitatore per mettere a proprio agio gli utenti in caso di stallo prolungato. Lo sperimentatore è rimasto invariato per tutti i test svolti.

\section{Scelta degli utenti}

Trattandosi di un sito destinato ad un'ampia varietà di persone (dai fruitori occasionali di visite guidate ai volontari), sono stati scelti 10 utenti eterogenei: 5 per il Gruppo A e 5 per il Gruppo B. 
Prima dell'inizio della sessione, ciascun partecipante ha compilato il \textbf{questionario anagrafico} (oltre ad aver accettato le clausole del consenso informato GDPR). 

La tabella seguente riassume la suddivisione:

\begin{table}[H]
\centering
\renewcommand{\arraystretch}{1.4}
\small
\begin{tabular}{|c|c|c|c|c|c|c|}
\hline
\textbf{ID} & \textbf{Età} & \textbf{Genere} & \textbf{Titolo di studio} & \textbf{Competenza informatica} & \textbf{Gruppo (Ordine)} \\
\hline
U01 & 22 & M & Laurea      & Alta  & Gruppo A (V1 $\to$ V2) \\ \hline
U02 & 20 & F & Diploma     & Media & Gruppo A (V1 $\to$ V2) \\ \hline
U03 & 38 & M & Laurea      & Alta  & Gruppo A (V1 $\to$ V2) \\ \hline
U04 & 26 & F & Laurea      & Media & Gruppo A (V1 $\to$ V2) \\ \hline
U05 & 54 & M & Diploma     & Bassa & Gruppo A (V1 $\to$ V2) \\ \hline
U06 & 23 & F & Laurea      & Alta  & Gruppo B (V2 $\to$ V1) \\ \hline
U07 & 28 & M & Laurea      & Media & Gruppo B (V2 $\to$ V1) \\ \hline
U08 & 45 & F & Laurea      & Media & Gruppo B (V2 $\to$ V1) \\ \hline
U09 & 33 & M & Diploma     & Bassa & Gruppo B (V2 $\to$ V1) \\ \hline
U10 & 21 & F & Laurea      & Alta  & Gruppo B (V2 $\to$ V1) \\ \hline
\end{tabular}
\caption{Panoramica anagrafica dei partecipanti al test di usabilità ($n=10$).}
\label{tab:partecipanti}
\end{table}

Maggiori dettagli demografici e relative distribuzioni grafiche saranno illustrati approfonditamente nel capitolo dell'analisi dei risultati (Capitolo 6).

\section{Scelta dei compiti e Moduli}

L'esperimento è basato su \textbf{11 compiti} (task) pensati per coprire il ciclo di vita completo della piattaforma: dalla ricerca di informazioni base, alla registrazione, per poi valutare le funzionalità transazionali (prenotazioni) e le interfacce riservate ai ruoli di volontario e amministratore. 
Per standardizzare lo svolgimento, oltre al questionario anagrafico iniziale, è stato predisposto un \textbf{modulo strutturato} (riportato di seguito) compilato fisicamente dallo sperimentatore per annotare click, tempi, errori e i commenti \textit{Think Aloud} degli utenti in tempo reale.

\vspace{0.5cm}

% T1
\noindent
\framebox[\linewidth]{
\begin{minipage}{0.96\linewidth}
\vspace{0.2cm}
\textbf{T1. Dire dove si trova la sede dell'organizzazione} \\[0.15cm]
\textit{Risposta attesa:} Via Branze 38, Brescia \\
\textbf{Tempo limite:} 1 min \\[0.2cm]
\hrule \vspace{0.2cm}
Tempo impiegato: \rule{2.5cm}{0.4pt} \hfill Click: \rule{2cm}{0.4pt} \hfill Errori: \rule{2cm}{0.4pt} \\[0.2cm]
Completato: \quad [\quad] Sì \qquad [\quad] No \qquad [\quad] Con Aiuto \\[0.2cm]
\hrule \vspace{0.2cm}
\textbf{Note Sperimentatore:} \\[1.2cm]
\textbf{Commenti Utente:} \\[1.2cm]
\end{minipage}} \vspace{0.5cm}

% T2
\noindent
\framebox[\linewidth]{
\begin{minipage}{0.96\linewidth}
\vspace{0.2cm}
\textbf{T2. Registrati come nuovo utente e accedi} \\[0.15cm]
\textit{Azione attesa:} Log-in completato con successo. \\
\textbf{Tempo limite:} 2 min \\[0.2cm]
\hrule \vspace{0.2cm}
Tempo impiegato: \rule{2.5cm}{0.4pt} \hfill Click: \rule{2cm}{0.4pt} \hfill Errori: \rule{2cm}{0.4pt} \\[0.2cm]
Completato: \quad [\quad] Sì \qquad [\quad] No \qquad [\quad] Con Aiuto \\[0.2cm]
\hrule \vspace{0.2cm}
\textbf{Note Sperimentatore:} \\[1.2cm]
\textbf{Commenti Utente:} \\[1.2cm]
\end{minipage}} \vspace{0.5cm}

% T3
\noindent
\framebox[\linewidth]{
\begin{minipage}{0.96\linewidth}
\vspace{0.2cm}
\textbf{T3. Prenota una visita guidata a "Castello di Sera" del X (seconda riga) per 2 persone} \\[0.15cm]
\textit{Azione attesa:} Prenotazione generata nel profilo. \\
\textbf{Tempo limite:} 2 min \\[0.2cm]
\hrule \vspace{0.2cm}
Tempo impiegato: \rule{2.5cm}{0.4pt} \hfill Click: \rule{2cm}{0.4pt} \hfill Errori: \rule{2cm}{0.4pt} \\[0.2cm]
Completato: \quad [\quad] Sì \qquad [\quad] No \qquad [\quad] Con Aiuto \\[0.2cm]
\hrule \vspace{0.2cm}
\textbf{Note Sperimentatore:} \\[1.2cm]
\textbf{Commenti Utente:} \\[1.2cm]
\end{minipage}} \vspace{0.5cm}

% T4
\noindent
\framebox[\linewidth]{
\begin{minipage}{0.96\linewidth}
\vspace{0.2cm}
\textbf{T4. Trova il codice della prenotazione appena effettuata} \\[0.15cm]
\textit{Risposta attesa:} Dettatura del token univoco generato randomicamente dal sistema. \\
\textbf{Tempo limite:} 2 min \\[0.2cm]
\hrule \vspace{0.2cm}
Tempo impiegato: \rule{2.5cm}{0.4pt} \hfill Click: \rule{2cm}{0.4pt} \hfill Errori: \rule{2cm}{0.4pt} \\[0.2cm]
Completato: \quad [\quad] Sì \qquad [\quad] No \qquad [\quad] Con Aiuto \\[0.2cm]
\hrule \vspace{0.2cm}
\textbf{Note Sperimentatore:} \\[1.2cm]
\textbf{Commenti Utente:} \\[1.2cm]
\end{minipage}} \vspace{0.5cm}

% T5
\noindent
\framebox[\linewidth]{
\begin{minipage}{0.96\linewidth}
\vspace{0.2cm}
\textbf{T5. Cancella la prenotazione appena effettuata} \\[0.15cm]
\textit{Azione attesa:} Modale di conferma d'annullamento superato. \\
\textbf{Tempo limite:} 1 min \\[0.2cm]
\hrule \vspace{0.2cm}
Tempo impiegato: \rule{2.5cm}{0.4pt} \hfill Click: \rule{2cm}{0.4pt} \hfill Errori: \rule{2cm}{0.4pt} \\[0.2cm]
Completato: \quad [\quad] Sì \qquad [\quad] No \qquad [\quad] Con Aiuto \\[0.2cm]
\hrule \vspace{0.2cm}
\textbf{Note Sperimentatore:} \\[1.2cm]
\textbf{Commenti Utente:} \\[1.2cm]
\end{minipage}} \vspace{0.5cm}

% T6
\noindent
\framebox[\linewidth]{
\begin{minipage}{0.96\linewidth}
\vspace{0.2cm}
\textbf{T6. Effettua il logout dall'utente attuale ed esegui il login come volontario (anna@example.com)} \\[0.15cm]
\textit{Azione attesa:} Switch dell'account superato con password: password123. \\
\textbf{Tempo limite:} 1 min \\[0.2cm]
\hrule \vspace{0.2cm}
Tempo impiegato: \rule{2.5cm}{0.4pt} \hfill Click: \rule{2cm}{0.4pt} \hfill Errori: \rule{2cm}{0.4pt} \\[0.2cm]
Completato: \quad [\quad] Sì \qquad [\quad] No \qquad [\quad] Con Aiuto \\[0.2cm]
\hrule \vspace{0.2cm}
\textbf{Note Sperimentatore:} \\[1.2cm]
\textbf{Commenti Utente:} \\[1.2cm]
\end{minipage}} \vspace{0.5cm}

% T7
\noindent
\framebox[\linewidth]{
\begin{minipage}{0.96\linewidth}
\vspace{0.2cm}
\textbf{T7. Imposta la disponibilità a un giorno qualsiasi non già impostato (es 25/04/26)} \\[0.15cm]
\textit{Azione attesa:} Slot sul calendario personale indicato come disp. \\
\textbf{Tempo limite:} 1 min \\[0.2cm]
\hrule \vspace{0.2cm}
Tempo impiegato: \rule{2.5cm}{0.4pt} \hfill Click: \rule{2cm}{0.4pt} \hfill Errori: \rule{2cm}{0.4pt} \\[0.2cm]
Completato: \quad [\quad] Sì \qquad [\quad] No \qquad [\quad] Con Aiuto \\[0.2cm]
\hrule \vspace{0.2cm}
\textbf{Note Sperimentatore:} \\[1.2cm]
\textbf{Commenti Utente:} \\[1.2cm]
\end{minipage}} \vspace{0.5cm}

% T8
\noindent
\framebox[\linewidth]{
\begin{minipage}{0.96\linewidth}
\vspace{0.2cm}
\textbf{T8. Effettua il logout ed esegui il login come amministratore (admin@example.com)} \\[0.15cm]
\textit{Azione attesa:} Dashboard admin raggiunta con successo (password123). \\
\textbf{Tempo limite:} 30 s \\[0.2cm]
\hrule \vspace{0.2cm}
Tempo impiegato: \rule{2.5cm}{0.4pt} \hfill Click: \rule{2cm}{0.4pt} \hfill Errori: \rule{2cm}{0.4pt} \\[0.2cm]
Completato: \quad [\quad] Sì \qquad [\quad] No \qquad [\quad] Con Aiuto \\[0.2cm]
\hrule \vspace{0.2cm}
\textbf{Note Sperimentatore:} \\[1.2cm]
\textbf{Commenti Utente:} \\[1.2cm]
\end{minipage}} \vspace{0.5cm}

% T9
\noindent
\framebox[\linewidth]{
\begin{minipage}{0.96\linewidth}
\vspace{0.2cm}
\textbf{T9. Crea una visita per la data impostata come disponibile da volontario prima (es. 25/04/26)} \\[0.15cm]
\textit{Azione attesa:} Creare l'evento associando la guida "Anna". \\
\textbf{Tempo limite:} 2 min \\[0.2cm]
\hrule \vspace{0.2cm}
Tempo impiegato: \rule{2.5cm}{0.4pt} \hfill Click: \rule{2cm}{0.4pt} \hfill Errori: \rule{2cm}{0.4pt} \\[0.2cm]
Completato: \quad [\quad] Sì \qquad [\quad] No \qquad [\quad] Con Aiuto \\[0.2cm]
\hrule \vspace{0.2cm}
\textbf{Note Sperimentatore:} \\[1.2cm]
\textbf{Commenti Utente:} \\[1.2cm]
\end{minipage}} \vspace{0.5cm}

% T10
\noindent
\framebox[\linewidth]{
\begin{minipage}{0.96\linewidth}
\vspace{0.2cm}
\textbf{T10. Elimina una visita da un luogo qualsiasi} \\[0.15cm]
\textit{Azione attesa:} Eliminazione riuscita e rimozione dall'elenco. \\
\textbf{Tempo limite:} 2 min \\[0.2cm]
\hrule \vspace{0.2cm}
Tempo impiegato: \rule{2.5cm}{0.4pt} \hfill Click: \rule{2cm}{0.4pt} \hfill Errori: \rule{2cm}{0.4pt} \\[0.2cm]
Completato: \quad [\quad] Sì \qquad [\quad] No \qquad [\quad] Con Aiuto \\[0.2cm]
\hrule \vspace{0.2cm}
\textbf{Note Sperimentatore:} \\[1.2cm]
\textbf{Commenti Utente:} \\[1.2cm]
\end{minipage}} \vspace{0.5cm}

% T11
\noindent
\framebox[\linewidth]{
\begin{minipage}{0.96\linewidth}
\vspace{0.2cm}
\textbf{T11. Cambia password da password123 a Admin123} \\[0.15cm]
\textit{Azione attesa:} Raggiungere le impostazioni del profilo Admin e mutare le credenziali. \\
\textbf{Tempo limite:} 1 min \\[0.2cm]
\hrule \vspace{0.2cm}
Tempo impiegato: \rule{2.5cm}{0.4pt} \hfill Click: \rule{2cm}{0.4pt} \hfill Errori: \rule{2cm}{0.4pt} \\[0.2cm]
Completato: \quad [\quad] Sì \qquad [\quad] No \qquad [\quad] Con Aiuto \\[0.2cm]
\hrule \vspace{0.2cm}
\textbf{Note Sperimentatore:} \\[1.2cm]
\textbf{Commenti Utente:} \\[1.2cm]
\end{minipage}} \vspace{0.5cm}
